\chapter{Kan Extensions}
\label{ch:kan}

\section{The Fundamental Question}

Consider functors $K: \cat{C} \to \cat{D}$ and $F: \cat{C} \to \cat{E}$. Can we ``extend'' $F$ through $K$ to obtain a functor $\cat{D} \to \cat{E}$? That is, can we fill in the dashed arrow:
\[
  \begin{tikzcd}
    \cat{C} \arrow[r, "K"] \arrow[dr, "F"'] & \cat{D} \arrow[d, dashed] \\
    & \cat{E}
  \end{tikzcd}
\]

In general, no extension exists on the nose. But we can ask for the \emph{best approximation}. There are two canonical choices: the \emph{left Kan extension} (which approximates from below, via colimits) and the \emph{right Kan extension} (which approximates from above, via limits). Kan extensions subsume virtually every other universal construction in category theory—limits, colimits, adjunctions, and more.

As Mac Lane famously wrote: ``The notion of Kan extensions subsumes all the other fundamental concepts of category theory.''

\section{Right Kan Extensions}

\begin{definition}[Right Kan extension]
Let $K: \cat{C} \to \cat{D}$ and $F: \cat{C} \to \cat{E}$ be functors. A \textbf{right Kan extension} of $F$ along $K$ is a pair $(\Ran_K F, \varepsilon)$ where:
\begin{itemize}
  \item $\Ran_K F: \cat{D} \to \cat{E}$ is a functor.
  \item $\varepsilon: (\Ran_K F) \circ K \Rightarrow F$ is a natural transformation (the \textbf{counit}).
\end{itemize}
satisfying the universal property: for any functor $G: \cat{D} \to \cat{E}$ and natural transformation $\alpha: GK \Rightarrow F$, there is a unique natural transformation $\bar\alpha: G \Rightarrow \Ran_K F$ such that $\varepsilon \circ (\bar\alpha K) = \alpha$.
\[
  \begin{tikzcd}
    \cat{C} \arrow[r, "K"] \arrow[dr, "F"'] & \cat{D} \arrow[d, "\Ran_K F"] \arrow[dl, "G"', dashed, bend right] \\
    & \cat{E}
  \end{tikzcd}
  \quad \varepsilon: (\Ran_K F)K \Rightarrow F, \quad \bar\alpha: G \Rightarrow \Ran_K F
\]
Equivalently, $\Ran_K F$ represents the functor
\[
  G \;\mapsto\; \Nat(GK, F),
\]
giving a natural bijection $\Nat(G, \Ran_K F) \cong \Nat(GK, F)$.
\end{definition}

\section{Left Kan Extensions}

\begin{definition}[Left Kan extension]
A \textbf{left Kan extension} of $F: \cat{C} \to \cat{E}$ along $K: \cat{C} \to \cat{D}$ is a pair $(\Lan_K F, \eta)$ where:
\begin{itemize}
  \item $\Lan_K F: \cat{D} \to \cat{E}$ is a functor.
  \item $\eta: F \Rightarrow (\Lan_K F) \circ K$ is a natural transformation (the \textbf{unit}).
\end{itemize}
satisfying the universal property: for any $G: \cat{D} \to \cat{E}$ and $\beta: F \Rightarrow GK$, there is a unique $\bar\beta: \Lan_K F \Rightarrow G$ such that $(\bar\beta K) \circ \eta = \beta$.

Equivalently, $\Lan_K F$ corepresents the functor $G \mapsto \Nat(F, GK)$:
\[
  \Nat(\Lan_K F, G) \;\cong\; \Nat(F, GK).
\]
\end{definition}

\begin{remark}[Duality]
Left and right Kan extensions are dual: $\Ran_K F$ in $\cat{E}$ equals $\Lan_K (F\op)$ in $\cat{E}\op$ (roughly). One can always reduce one to the other by dualizing.
\end{remark}

\section{Pointwise Formula for Kan Extensions}

When $\cat{E}$ is complete (resp.\ cocomplete), Kan extensions exist and can be computed by an explicit (co)limit formula.

\begin{theorem}[Pointwise right Kan extension]
\label{thm:pointwise-ran}
Let $\cat{C}$ be small, $K: \cat{C} \to \cat{D}$, $F: \cat{C} \to \cat{E}$. If $\cat{E}$ has all small limits, then $\Ran_K F$ exists and is given, for $d \in \cat{D}$, by:
\[
  (\Ran_K F)(d) = \lim\bigl(K \downarrow d \xrightarrow{\pi} \cat{C} \xrightarrow{F} \cat{E}\bigr),
\]
where $K \downarrow d$ is the comma category (objects: pairs $(c, f: Kc \to d)$, morphisms: maps in $\cat{C}$ over $d$) and $\pi$ is the projection functor.
\end{theorem}

\begin{theorem}[Pointwise left Kan extension]
\label{thm:pointwise-lan}
Under dual hypotheses ($\cat{E}$ cocomplete):
\[
  (\Lan_K F)(d) = \colim\bigl(d \downarrow K \xrightarrow{\pi} \cat{C} \xrightarrow{F} \cat{E}\bigr),
\]
where $d \downarrow K$ has objects $(c, f: d \to Kc)$.
\end{theorem}

\begin{proof}[Proof of Theorem~\ref{thm:pointwise-lan}]
We verify the universal property. Given $G: \cat{D} \to \cat{E}$ and $\beta: F \Rightarrow GK$, we need a unique $\bar\beta: \Lan_K F \Rightarrow G$.

At a fixed $d \in \cat{D}$, we need a map $(\Lan_K F)(d) = \colim_{(c,f) \in d\downarrow K} F(c) \to G(d)$. For each $(c, f: d \to Kc)$, the component $\beta_c: F(c) \to GKc$, composed with $G(f): GKc \to Gd$, gives $G(f) \circ \beta_c: F(c) \to Gd$. These are compatible with morphisms in $d \downarrow K$, hence define a cocone, hence a unique map from the colimit. Naturality in $d$ is verified by functoriality of everything involved.

Uniqueness: the unit $\eta_c: F(c) \to (\Lan_K F)(Kc)$ corresponds to the map from $F(c)$ to the colimit indexed by $(c, \id_{Kc}) \in Kc \downarrow K$. Any $\bar\beta: \Lan_K F \Rightarrow G$ with $(\bar\beta K) \circ \eta = \beta$ determines, for each $(c, f: d \to Kc)$, the map $F(c) \xrightarrow{\eta_c} (\Lan_K F)(Kc) \xrightarrow{(\Lan_K F)(f)} (\Lan_K F)(d) \xrightarrow{\bar\beta_d} Gd$; since the colimit is generated by the elements of $F(c)$ over all $(c,f)$, $\bar\beta_d$ is determined.
\end{proof}

\section{Adjunctions as Kan Extensions}

\begin{theorem}
\label{thm:adj-as-kan}
Let $F: \cat{C} \to \cat{D}$ and $G: \cat{D} \to \cat{C}$. Then $F \dashv G$ if and only if $G \cong \Ran_F \Id_{\cat{C}}$ (and the Kan extension is pointwise), or equivalently, $F \cong \Lan_G \Id_{\cat{D}}$.
\end{theorem}

\begin{proof}
$\Nat(G', \Ran_F \Id_{\cat{C}}) \cong \Nat(G'F, \Id_{\cat{C}}) = $ \{natural transformations $G'F \Rightarrow \Id_{\cat{C}}\}$. This is the universal property of the right adjoint: $G$ represents $\Nat(-F, \Id_{\cat{C}})$, i.e., $F \dashv G$ iff $G = \Ran_F \Id_{\cat{C}}$.
\end{proof}

\begin{corollary}
If $F \dashv G$, then for any $H: \cat{C} \to \cat{E}$, $\Ran_F H \cong H \circ G$ (and the extension is absolute, computed without colimits).
\end{corollary}

\begin{proof}
$\Nat(L, \Ran_F H) \cong \Nat(LF, H) \cong \Nat(L, HG)$, so $\Ran_F H \cong HG$.
\end{proof}

\section{The Density Theorem}

The density theorem is a fundamental result relating any functor to its own Kan extension along itself.

\begin{theorem}[Density theorem]
\label{thm:density}
Let $K: \cat{C} \to \cat{D}$ be any functor. Then the left Kan extension of $K$ along itself is (naturally isomorphic to) the identity:
\[
  \Lan_K K \;\cong\; \Id_{\cat{D}}.
\]
More precisely, if $\cat{D}$ is cocomplete and $K$ is small (e.g., $\cat{C}$ is small), then the left Kan extension $\Lan_K K: \cat{D} \to \cat{D}$ is isomorphic to $\Id_{\cat{D}}$, i.e., every object of $\cat{D}$ is a colimit of objects in the image of $K$.
\end{theorem}

\begin{proof}
Using the pointwise formula: $(\Lan_K K)(d) = \colim_{(c, Kc \to d)} Kc$. The colimit is taken over the comma category $(d \downarrow K)\op$, which is the category of objects in the image of $K$ with maps to $d$. The universal property of this colimit gives the identity morphism on $d$.

More abstractly: by Theorem~\ref{thm:adj-as-kan}, $\Lan_K K \cong \Id_\cat{D}$ iff the pair $(\Id_\cat{D}, K)$ constitutes an adjunction—but actually the density theorem is stronger and doesn't require an adjunction. The key point is that every $d \in \cat{D}$ is canonically a colimit of objects $Kc$ via the comma category $d \downarrow K$.
\end{proof}

\begin{corollary}[Yoneda as a density statement]
Every presheaf $F \in [\cat{C}\op, \Set]$ is a canonical colimit of representables:
\[
  F \cong \colim_{(c, \yo(c) \to F)} \yo(c).
\]
(The colimit is over the category of elements of $F$.) This is the density theorem for the Yoneda embedding $\yo: \cat{C} \to [\cat{C}\op, \Set]$.
\end{corollary}

\section{Nerve and Geometric Realization}

A classic example of Kan extensions in topology.

\begin{example}[Nerve and realization]
Let $\Delta$ be the simplex category: objects are the finite ordinals $[n] = \{0,1,\ldots,n\}$ and morphisms are weakly order-preserving maps. The standard simplices form a functor $\Delta^\bullet: \Delta \to \Top$, $[n] \mapsto \Delta^n$ (the standard topological $n$-simplex).

The \textbf{singular nerve} of a topological space $X$ is the presheaf $N(X) = \Top(\Delta^\bullet, X): \Delta\op \to \Set$, i.e., $N(X)_n = \Top(\Delta^n, X)$ (singular $n$-simplices in $X$).

The \textbf{geometric realization} of a simplicial set $S: \Delta\op \to \Set$ is:
\[
  |S| = \Lan_{\yo} (\Delta^\bullet)(S) = \int^{[n]} S_n \times \Delta^n \quad (\text{a coend/colimit formula}).
\]
This is the left Kan extension of $\Delta^\bullet: \Delta \to \Top$ along the Yoneda embedding $\yo: \Delta \to [\Delta\op, \Set]$:
\[
  \begin{tikzcd}
    \Delta \arrow[r, "\yo"] \arrow[dr, "\Delta^\bullet"'] & {[\Delta\op, \Set]} \arrow[d, "|-|"] \\
    & \Top
  \end{tikzcd}
\]
The adjunction $|-| \dashv N$ (geometric realization left adjoint to nerve) is the Kan extension adjunction.
\end{example}

\section{Day Convolution}

\begin{theorem}[Day convolution]
Let $(\cat{C}, \otimes, I)$ be a small symmetric monoidal category. The presheaf category $[\cat{C}\op, \Set]$ carries a symmetric monoidal structure given by \textbf{Day convolution}:
\[
  (F * G)(c) = \int^{a,b \in \cat{C}} \cat{C}(c, a \otimes b) \times Fa \times Gb \;=\; \colim_{c \to a \otimes b} Fa \times Gb,
\]
(a coend/left Kan extension formula). This is the left Kan extension of the monoidal product $\otimes: \cat{C} \times \cat{C} \to \cat{C}$ along $\yo \times \yo: \cat{C} \times \cat{C} \to [\cat{C}\op,\Set] \times [\cat{C}\op,\Set]$, composed with $\yo$.

The Yoneda embedding $\yo: \cat{C} \hookrightarrow [\cat{C}\op,\Set]$ is strong monoidal: $\yo(a) * \yo(b) \cong \yo(a \otimes b)$.
\end{theorem}

Day convolution unifies many constructions: formal power series (polynomial rings), Dirichlet series, convolution algebras of functions, and the tensor product of $\mathbf{Ab}$-enriched presheaves.

\section{Global Kan Extension Functors}

\begin{theorem}
Let $K: \cat{C} \to \cat{D}$ be a functor. The precomposition functor
\[
  K^*: [\cat{D}, \cat{E}] \to [\cat{C}, \cat{E}], \qquad G \mapsto G \circ K,
\]
has a left adjoint $\Lan_K: [\cat{C},\cat{E}] \to [\cat{D},\cat{E}]$ (global left Kan extension) and a right adjoint $\Ran_K: [\cat{C},\cat{E}] \to [\cat{D},\cat{E}]$ (global right Kan extension), whenever $\cat{E}$ is cocomplete and complete respectively.

Symbolically: $\Lan_K \dashv K^* \dashv \Ran_K$.
\end{theorem}

\begin{proof}
For $\Lan_K$: the natural bijection $[\cat{D},\cat{E}]({\Lan_K F}, G) \cong [\cat{C},\cat{E}](F, K^* G) = \Nat(F, GK)$ is the defining universal property of $\Lan_K F$. For $\Ran_K$: similar, with $\Nat(K^*G, F) = \Nat(GK, F) \cong [\cat{D},\cat{E}](G, \Ran_K F)$.
\end{proof}

\begin{corollary}
When $K = \yo: \cat{C} \to [\cat{C}\op,\Set]$ is the Yoneda embedding, $\Lan_\yo: [[\cat{C}\op,\Set], \cat{E}] \to [\cat{C},\cat{E}]$ realizes the free cocompletion universal property: every functor $F: \cat{C} \to \cat{E}$ (with $\cat{E}$ cocomplete) extends uniquely to a cocontinuous $\bar F: [\cat{C}\op,\Set] \to \cat{E}$, and this extension is $\Lan_\yo F$.
\end{corollary}

\section{Absolute Kan Extensions}

\begin{definition}[Absolute Kan extension]
A left Kan extension $(\Lan_K F, \eta)$ is \textbf{absolute} if it is preserved by all functors $H: \cat{E} \to \cat{E}'$: i.e., $H \circ \Lan_K F \cong \Lan_K (H \circ F)$ (with the obvious unit).
\end{definition}

\begin{example}
If $F \dashv G$, then $\Ran_F \Id_{\cat{C}} \cong G$ (Theorem~\ref{thm:adj-as-kan}), and this Kan extension is absolute. Absolute Kan extensions are precisely those computed without using (co)limits—they are ``preserved for free'' by all functors.
\end{example}

\begin{theorem}
A Kan extension is absolute if and only if it is computed by the counit (resp.\ unit) of an adjunction.
\end{theorem}

\section{Kan Extensions in 2-Categories}

The notion of Kan extension generalizes to any 2-category.

\begin{definition}[Kan extension in a 2-category]
In a 2-category $\mathcal{K}$, given 1-cells $K: C \to D$ and $F: C \to E$, a \textbf{left Kan extension} is a pair $(\Lan_K F: D \to E, \eta: F \Rightarrow (\Lan_K F) \circ K)$ satisfying the obvious 2-categorical universal property.
\end{definition}

In $\Cat$, this recovers the usual Kan extensions. In other 2-categories (e.g., the 2-category of rings, bimodules, bimodule maps), it gives Morita-theoretic Kan extensions.

\section{Coends and the Co-Yoneda Lemma}

Kan extensions are intimately connected with coends, a useful computational tool.

\begin{definition}[Coend]
Let $T: \cat{C}\op \times \cat{C} \to \cat{E}$ be a functor (a \textbf{dinatural} situation). The \textbf{coend} $\int^c T(c,c)$ is the colimit of the diagram indexed by the \textbf{twisted arrow category} of $\cat{C}$: objects are morphisms $f: c \to c'$ in $\cat{C}$, with maps $T(f, \id_c): T(c',c) \to T(c,c)$ and $T(\id_{c'}, f): T(c',c) \to T(c',c')$.

Equivalently, $\int^c T(c,c) = \colim_{f: c \to c' \text{ in } \cat{C}} \bigl(T(c',c) \rightrightarrows T(c,c) \sqcup T(c',c')\bigr)$ (coequalizer of action on morphisms).
\end{definition}

\begin{theorem}[Co-Yoneda lemma]
For $F: \cat{C} \to \cat{E}$ and $H \in [\cat{C}\op, \Set]$:
\[
  \int^c Hc \cdot Fc \;\cong\; \Lan_\yo F (H),
\]
where $Hc \cdot Fc$ denotes the \textbf{tensor} (copower): $\bigsqcup_{h \in Hc} Fc$ (or $Hc \times Fc$ in $\Set$).

In particular, for $H = \yo(d) = \cat{C}(-, d)$:
\[
  \int^c \cat{C}(c, d) \cdot Fc \;\cong\; Fd.
\]
\end{theorem}

This ``co-Yoneda'' formula is dual to the Yoneda lemma and shows that any functor value $Fd$ can be reconstructed as a coend.

\section{Exercises}

\begin{exercise}
Show that the right Kan extension along the unique functor $!: \cat{C} \to \mathbf{1}$ is the limit: $\Ran_! F \cong \lim F$. Dualize for left Kan extensions and colimits.
\end{exercise}

\begin{exercise}
Verify that $\Ran_{\Id} F \cong F$ (the right Kan extension along the identity is $F$ itself).
\end{exercise}

\begin{exercise}
Show that if $K: \cat{C} \to \cat{D}$ is fully faithful, then the unit $\eta: F \Rightarrow (\Lan_K F) \circ K$ is a natural isomorphism (the Kan extension truly extends $F$ in this case).
\end{exercise}

\begin{exercise}
Use the pointwise formula to compute $\Lan_K F$ when $K: \{*\} \to \cat{D}$ picks out a single object $d \in \cat{D}$ and $F: \{*\} \to \cat{E}$ picks out $e \in \cat{E}$. What is $(\Lan_K F)(d')$ for varying $d' \in \cat{D}$?
\end{exercise}

\begin{exercise}[$\star$]
Prove the density theorem (Theorem~\ref{thm:density}) for the Yoneda embedding: every presheaf $F \in [\cat{C}\op,\Set]$ is a colimit of representables, with the colimit indexed by the category of elements of $F$.
\end{exercise}

\begin{exercise}[$\star$]
Using Day convolution, describe the monoidal structure on $[\mathbb{N}, \Set]$ (presheaves on the monoid $(\mathbb{N}, +, 0)$ viewed as a one-object category). How does this relate to formal power series?
\end{exercise}

\begin{exercise}[$\star\star$]
Prove that for any adjunction $F \dashv G$, the right Kan extension $\Ran_F H \cong HG$ for any $H: \cat{D} \to \cat{E}$. (This is the ``absolute'' Kan extension claim.) What is the counit of this Kan extension?
\end{exercise}

\begin{exercise}[$\star\star$]
Let $i: \cat{C} \hookrightarrow \cat{D}$ be a full subcategory inclusion with $\cat{D}$ cocomplete. Show that the left Kan extension $\Lan_i F$ for $F: \cat{C} \to \cat{E}$ gives the ``best approximation'' to $F$ on $\cat{D}$, and that if $F$ has a left adjoint on $\cat{C}$, the Kan extension has a left adjoint on $\cat{D}$ (under mild conditions).
\end{exercise}
